% !TeX root = ../组合数学笔记.tex
% 由主文档通过 \input 载入；请编译主文档。
\setcounter{chapter}{1} % 保留原书章号。
\chapter{排列、组合与分配模型}
\label{ch:elementary-counting}

排列与组合公式处理的是几种标准对象。使用前先回答四个问题：元素是否可区分，顺序是否保留，是否允许重复，是否要求每个容器非空。把限制直接写进对象定义，比从题目中的“选”“排”“分”猜测公式更可靠。

\section{排列、组合与可重复选择}

\begin{definition}
从 $n$ 个不同元素中选出 $k$ 个不同元素并排成序列，称为一个 $k$-\emph{排列（permutation）}；只保留所选元素的集合，称为一个 $k$-\emph{组合（combination）}。当 $k=n$ 时，排列也称全排列。
\end{definition}

对 $0\le k\le n$，逐位选择得到排列数 $(n)_k=n(n-1)\cdots(n-k+1)=n!/(n-k)!$，其中 $(n)_0=1$。每个 $k$ 元子集恰有 $k!$ 个排列，因此
\begin{equation}
  \binom nk=\frac{(n)_k}{k!}=\frac{n!}{k!(n-k)!}.
  \label{eq:binomial-count}
\end{equation}
若允许每位重复选取，则长度为 $k$ 的序列是从位置集 $[k]$ 到元素集 $[n]$ 的映射，共有 $n^k$ 个。空位置集只有一个映射，故在这种计数语境中约定 $0^0=1$；非空集合到空集没有映射。

\begin{example}
用十个数字组成四位十进制正整数。允许重复时，首位有 $9$ 种，其余各有 $10$ 种，得 $9000$ 个；数字互不相同时，依次有 $9,9,8,7$ 种选择，得 $4536$ 个。第二位仍有 $9$ 种，是因为可以用零，但不能重复首位。这说明分步选择数可以相同，而对应的候选集合并不相同。
\end{example}

\section{多重集合排列与圆排列}

一个\emph{多重集合（multiset）}允许元素具有重数。若 $r$ 类元素的重数为 $n_1,\ldots,n_r\ge0$，总数为 $n$，先把所有副本临时标为不同，得到 $n!$ 个排列。去掉副本标签后，每个实际排列恰有 $n_1!\cdots n_r!$ 个原像，故排列数为
\begin{equation}
  \binom{n}{n_1,\ldots,n_r}=\frac{n!}{n_1!\cdots n_r!}.
  \label{eq:multinomial-count}
\end{equation}
例如字母串 $\mathrm{BANANA}$ 有三个 $A$、两个 $N$、一个 $B$，不同排列数为 $6!/(3!2!)=60$。

把 $n\ge1$ 个不同的人围成圆桌，只把整体旋转视为相同。固定一人的位置，其余 $n-1$ 人按顺时针方向排列，得到 $(n-1)!$ 种。固定位置是在每个等价类中选唯一代表；若再把镜像视为相同，等价关系已改变。含重复颜色时也不能一般地用线性排列数除以 $n$：例如二色交替排列可能在未转满一圈前恢复原状，各轨道未必等大。

\begin{example}
五个不同的人围成圆桌，其中甲乙必须相邻。把甲乙合成一个块后，有四个不同的圆排列单元，共 $(4-1)!$ 种，块内有 $2$ 种顺序，故共有 $12$ 种。不要求相邻时有 $24$ 种，因此甲乙不相邻也有 $12$ 种。此处块与圆桌安排可相互恢复，计数没有遗漏或重复。
\end{example}

\section{隔板法与整数向量}

将 $n\ge0$ 个相同物品放入 $k\ge1$ 个不同盒子，由各盒数量 $(x_1,\ldots,x_k)$ 唯一确定方案。它等价于非负整数方程 $x_1+\cdots+x_k=n$。

\begin{theorem}[隔板法（stars and bars）]
上述方程的非负整数解数为 $\binom{n+k-1}{k-1}$。若各变量均为正整数，则在 $n\ge k$ 时解数为 $\binom{n-1}{k-1}$，在 $n<k$ 时为零。
\end{theorem}
\begin{proof}
把 $n$ 个星号排成一行，插入 $k-1$ 个隔板，按隔板划分成 $k$ 段，段内星号数即为 $x_i$；允许隔板相邻和出现在两端。由此与含 $n$ 个星号、$k-1$ 个隔板的字符串建立双射，选隔板位置得到第一式。若 $x_i\ge1$，令 $y_i=x_i-1\ge0$，总和变成 $n-k$，应用第一式即得第二式。
\end{proof}

更一般地，若 $x_i\ge a_i$，其中 $a_i$ 为给定非负整数，先令 $y_i=x_i-a_i$；剩余总数 $n-\sum_i a_i$ 为负时无解。若同时有上界，则隔板位置受到额外约束，不能继续直接套用无约束公式。

\begin{example}
方程 $x_1+x_2+x_3=8$ 有 $\binom{10}{2}=45$ 个非负整数解；各变量均正时有 $\binom72=21$ 个；若 $x_1\ge2,x_2\ge1,x_3\ge0$，平移后总和为 $5$，有 $\binom72=21$ 个。后两者数量相同，但对应的原始解集不同。
\end{example}

\section{分配问题的统一分类}

不同物品放入不同盒子是映射问题：不限制空盒时为任意映射，每盒至多一件时为单射，每盒至少一件时为满射。相同物品放入不同盒子则是整数向量问题。若不同物品的盒子没有名称且只保留非空盒，得到\emph{集合分拆（set partition）}；若物品与盒子都相同，只保留各盒的正数量并按非增次序记录，得到\emph{整数分拆（integer partition）}。后两者将在第\ref{ch:egf}、\ref{ch:classical-structures}章讨论。

把正整数 $n$ 写成有序正整数之和称为\emph{整数合成（composition）}。固定 $k$ 项时有 $\binom{n-1}{k-1}$ 种；不固定项数而 $n\ge1$ 时，在一排 $n$ 个单位的 $n-1$ 个间隙中任意选择断点，故共有 $2^{n-1}$ 种。$1+3$ 与 $3+1$ 是不同合成，却是同一个整数分拆。是否保留顺序决定了模型。

\section{习题}
\begin{enumerate}
  \item\label{ex:elementary-counting-1} 求字母串 $\mathrm{MISSISSIPPI}$ 的不同排列数，并解释分母各因子的意义。
  \item\label{ex:elementary-counting-2} 六人排成一行，指定两人不相邻，有多少种排列？改为圆桌且只识别旋转后再求解。
  \item\label{ex:elementary-counting-3} 求 $x_1+x_2+x_3+x_4=12$ 且各 $x_i\ge2$ 的整数解数。
  \item\label{ex:elementary-counting-4} 把五个不同物品放入三个不同盒子，分别写出任意分配、每盒至多一件、每盒至少一件的模型；最后一种暂不要求套公式。
\end{enumerate}

